\section{A sharp Picard bound for toric bundles}\label{sec:interactions}

We now consider bundles with fibre a product of degree-six del Pezzo
surfaces and base a product of projective spaces. In this class the
maximal Picard number of a stable example is $\lfloor5n/4\rfloor+1$
for every $n\ge4$. The upper bound allows arbitrary integral twists.
Its proof combines three requirements: Fano positivity bounds the
twisting at each base factor, stability requires each surface factor
to meet at least two base factors, and stability excludes a product
decomposition. The constructions in the following sections show that
these requirements can be attained simultaneously.

\subsection{Bundle fans and a relative tangent obstruction}
Let $F$ be a smooth projective toric variety with lattice $N_F$, and let
$B=\prod_{c=1}^k\PP^{b_c}$, where $k,b_c\ge1$. Throughout this section a
\emph{toric bundle with fibre $F$ over $B$} has the following explicit
fan presentation. Its lattice is $N_F\oplus\bigoplus_c\ZZ^{b_c}$;
its fibre rays are those of $F$ in $N_F$; and its other primitive rays are
\begin{equation}\label{interact:rays}
 f^c_a=(0,e^c_a)\quad(1\le a\le b_c),\qquad
 f^c_0=\left(t_c,-\sum_{a=1}^{b_c}e^c_a\right),\qquad t_c\in N_F.
\end{equation}
For every maximal fibre cone and every choice of one omitted ray from
each set $\{f^c_0,\ldots,f^c_{b_c}\}$, their remaining generators form
a maximal cone of $X$, and these are all the maximal cones.
This is the fan description of a locally trivial toric bundle
\cite[Theorem~3.3.19]{CoxLittleSchenck2011}; it is part of our hypotheses,
and no assertion about an arbitrary Mori fibre space is intended.
The vectors $t_c$ specify the twisting. The displayed normalization is
obtained by an integral shear from the lifts of the standard base rays.

Write $P_F=\{x:\langle x,v\rangle\ge-1\text{ for every fibre ray }v\}$
when $F$ is Fano. All volume measures below are lattice normalized,
without a factorial. The first obstruction only requires that $P_F$
have barycentre zero. For the special case $F=\PP^1$ and base dimension
at least two, the strict inequality below recovers
\cite[Theorem~8.1]{BiswasDeyGencPoddar2021}.

\begin{proposition}[Obstruction over projective space]\label{interact:obstruction}
Let $F$ be a smooth toric Fano variety of positive dimension $f$, and let
$\pi:X\longrightarrow\PP^b$ be a toric bundle with the fan presentation~\eqref{interact:rays} and
fibre $F$, where $b\ge1$. Assume that $X$ is Fano and that the barycentre
of the anticanonical polytope $P_F$ is zero. Then
\[
 \mu_{-K_X}(T_{X/\PP^b})\ge\mu_{-K_X}(T_X).
\]
In particular, $T_X$ is not anticanonically stable. The inequality is
strict if $b\ge2$ and the twisting vector $t_1$ in~\eqref{interact:rays}
is nonzero. It is an equality if $b=1$ or $t_1=0$.
\end{proposition}
\begin{proof}
Write $n=f+b$ and specialize~\eqref{interact:rays} to a single base
factor. The lifted base rays are
\[
 (0,e_1),\ldots,(0,e_b),\qquad
 (t,-e_1-\cdots-e_b),\qquad t\in N_F.
\]
Put $s(x)=\langle x,t\rangle$ and $L(x)=b+1+s(x)$ for $x\in P_F$.
The anticanonical polytope of $X$ is
\[
 P_X=\{(x,y):x\in P_F,\ y_i\ge-1,
                   \ \sum_{i=1}^b(y_i+1)\le L(x)\}.
\]
Ampleness of $-K_X$ implies $L(x)>0$ on $P_F$. Indeed, the vertex of
$P_X$ corresponding to any maximal fibre cone together with the base
cone generated by $e_1,\ldots,e_b$ must satisfy the remaining base
inequality strictly. This gives $L(x)>0$ at every vertex of $P_F$, hence
throughout $P_F$ by affinity.

Let $A_q=\int_{P_F}L(x)^q\,dx$, with lattice-normalized Lebesgue measure.
The slices are $b$-simplices, so
\[
 D:=(-K_X)^n=\frac{n!}{b!}A_b.
\]
Each of the $b+1$ base-ray divisors has the same degree
\[
 d=\frac{(n-1)!}{(b-1)!}A_{b-1}.
\]
For the facets $y_i=-1$ this follows by integrating the corresponding
$(b-1)$-simplex volumes. All $b+1$ degrees are equal by the principal
divisor relations from the characters dual to $e_i$.
The relative tangent bundle has rank $f$ and first Chern class equal
to the sum of the fibre-ray divisors. Hence
\begin{align*}
 \mu(T_{X/\PP^b})-\mu(T_X)
 &=\frac{D-(b+1)d}{f}-\frac{D}{n}\\
 &=\frac{(n-1)!}{f(b-1)!}
   \left(A_b-(b+1)A_{b-1}\right)\\
 &=\frac{(n-1)!}{f(b-1)!}
   \int_{P_F}s(x)(b+1+s(x))^{b-1}\,dx.
\end{align*}
The barycentre assumption implies $\int_{P_F}s(x)\,dx=0$. Subtracting
$(b+1)^{b-1}\int s$ from the last integral rewrites it as
\[
 \int_{P_F}s(x)
 \left((b+1+s(x))^{b-1}-(b+1)^{b-1}\right)\,dx.
\]
This integrand is nonnegative, because the function $z\mapsto z^{b-1}$
is nondecreasing for $z>0$. If $b\ge2$ and $t\ne0$, it is strictly
positive whenever $s(x)\ne0$. Since $P_F$ is full-dimensional, that set
has positive measure. If $b=1$ or $t=0$, the integrand vanishes.
\end{proof}

\begin{corollary}[The Assarf--Nill extremal bundles]\label{interact:ANexamples}
Let $X_m$ be the toric variety associated with the smooth Fano polytope
$B_m$ of Assarf--Nill \cite[Definition~41 and Proposition~43]{AssarfNill2016}.
Thus $\dim X_m=3m$, $\rho(X_m)=4m+1$, and $X_m$ is a toric
$S_6^m$-bundle over $\PP^m$, with nonzero twist having one ray component
in each hexagon factor. Then $T_{X_m}$ is not anticanonically stable
for every $m\ge1$, and is unstable for every $m\ge2$.
\end{corollary}
\begin{proof}
The anticanonical polytope of $S_6^m$ is centrally symmetric and thus
has barycentre zero. Apply the proposition with $b=m$ and $f=2m$.
\end{proof}

\subsection{Incidences of hexagon fibre factors}
Let $S_6$ be the blow-up of $\PP^2$ at its three torus-fixed points.
Its fan rays, in a suitable lattice basis $r,q$, are
$\pm r,\pm q,\pm(r+q)$. Its anticanonical polytope is centrally symmetric.
For fibre $S_6^j$, decompose $N_F=\bigoplus_{i=1}^jN_i$ and write
$t_c=(t_{1c},\ldots,t_{jc})$. The \emph{incidence graph of the twisting}
is the bipartite graph on the fibre factors and base factors, with an
edge $ic$ precisely when $t_{ic}\ne0$. Thus its edges record which base
factors twist a given surface factor.

\begin{theorem}[Picard bound for hexagon bundles]\label{interact:bound}
Let $j,k\ge1$, let $b_1,\ldots,b_k\ge1$, and let $X$ be a locally
trivial toric $S_6^j$-bundle over $\prod_{c=1}^k\PP^{b_c}$, with the fan presentation~\eqref{interact:rays}. Assume that
$X$ is smooth Fano and that $T_X$ is anticanonically stable. Then
\[
 \rho(X)\le n+\lfloor n/4\rfloor+1,
 \qquad n=\dim X.
\]
Here the twists in each hexagon lattice are arbitrary integral vectors;
they are not required to have the graph form of Section~\ref{sec:graphs}.
More precisely, if $s=\sum_cb_c$ and $\ell$ is the number of edges of
the incidence graph, then
\begin{equation}\label{interact:incidence-inequalities}
 n=2j+s,\qquad \rho(X)=4j+k,\qquad
 2j\le\ell\le s,\qquad j+k-1\le\ell.
\end{equation}
\end{theorem}
\begin{proof}
Use the lattice decomposition and incidence graph defined above.
Let $\ell$ be the number of incidences and put $s=\sum_cb_c$.
Counting rays gives
\[
 n=2j+s,\qquad\rho(X)=4j+k.
\]

For each hexagon let $\psi_i$ be the piecewise linear function equal
to one on its primitive rays and linear on its fan cones. For an
integral nonzero vector $t$, the positive integer $\psi_i(t)$ is the
sum of its coefficients in the minimal containing unimodular cone.
The rays of the $c$th base factor form a primitive collection:
any proper subset lies in a lifted base cone, while the whole set does not.
Their sum is $t_c$ in the fibre lattice. Expressing it in the minimal
containing product cone gives a primitive relation of degree
\[
 b_c+1-\sum_i\psi_i(t_{ic}).
\]
The Fano condition makes this degree a positive integer
\cite[Theorem~6.4.9(b)]{CoxLittleSchenck2011}. Therefore
\[
 \#\{i:t_{ic}\ne0\}\le\sum_i\psi_i(t_{ic})\le b_c,
 \qquad \ell\le s.
\]

We next show that each hexagon vertex of the incidence graph has
degree at least two. Let $Q_i$ be its centrally symmetric
anticanonical hexagon, and define
\[
 L_c(x)=b_c+1+\sum_i\langle x_i,t_{ic}\rangle,\qquad
 G(x)=\prod_{c=1}^k\frac{L_c(x)^{b_c}}{b_c!},\qquad
 V=\int_{\prod_iQ_i}G(x)\,dx.
\]
The Fano assumption gives $L_c>0$ on the product of hexagons, by the
same strict vertex-inequality argument as in the preceding proposition.
The anticanonical polytope is a product of simplex slices over
$\prod_iQ_i$, so $V$ is its lattice volume.
Since the base has positive dimension, $N_{i,\CC}$ is a proper
subspace, and no base ray lies in it. Its normalized ray-weight sum is
\[
 w(N_{i,\CC})
 =\frac1V\sum_{H\text{ facet of }Q_i}
   \int_{H\times\prod_{h\ne i}Q_h}G(x)\,d\sigma\,dx_{-i}
 =2+\frac1V\int_{\prod_iQ_i}x_i\cdot\nabla_iG(x)\,dx.
\]
The second equality is the divergence theorem in the $x_i$ variables.
This is the slice-volume identity used in subspace concentration
\cite[Eq.~(4.4)]{HenkLinke2014}; here we use the symmetry of the fibre
factor, without assuming that the total polytope is centred.
For a facet with primitive outward normal $u$ and equation
$\langle u,x_i\rangle=1$, the normal component of $x_i$ is
$1/\|u\|$, exactly converting Euclidean facet measure to lattice
facet measure. Thus the identity has no missing normalization factor.

If all $t_{ic}$ vanish, then $w(N_{i,\CC})=2$ and stability fails.
If only $t_{ic}$ is nonzero, write
$z=\langle x_i,t_{ic}\rangle$ and
$A=b_c+1+\sum_{h\ne i}\langle x_h,t_{hc}\rangle>0$. Apart from a
positive factor independent of $x_i$, the inner integral contributing
to $\int x_i\cdot\nabla_iG$ is
\[
 b_c\int_{Q_i}z(A+z)^{b_c-1}\,dx_i
 =b_c\int_{Q_i}z\big((A+z)^{b_c-1}-A^{b_c-1}\big)\,dx_i\ge0.
\]
Here $\int_{Q_i}z=0$ by central symmetry, and positivity follows from
monotonicity as before. Consequently $w(N_{i,\CC})\ge2$, again
contrary to stability. Thus $\ell\ge2j$.

Finally, the incidence graph is connected. Otherwise, grouping the
lattice summands according to its components gives a direct sum of
proper subspaces containing all ray generators, so the ray matroid is
disconnected, contrary to Lemma~\ref{lem:raycomponents}. A connected graph on $j+k$
vertices has at least $j+k-1$ edges. Combining the inequalities gives
\[
 s\ge\ell\ge2j,\qquad s\ge\ell\ge j+k-1.
\]
Hence $n\ge4j$ and
\[
 \rho(X)=4j+k\le 4j+s-j+1=n+j+1
 \le n+\lfloor n/4\rfloor+1,
\]
as claimed.
\end{proof}

\begin{corollary}[Necessary conditions for equality]\label{interact:equality}
Under the hypotheses of Theorem~\ref{interact:bound}, suppose that
$n=4j+r$ with $0\le r\le3$ and that
$\rho(X)=\lfloor5n/4\rfloor+1=5j+r+1$. Then the incidence graph is a tree,
\[
 k=j+r+1,\qquad \ell=\sum_cb_c=2j+r,\qquad
 \sum_{i=1}^j\bigl(\deg(i)-2\bigr)=r.
\]
Every nonzero $t_{ic}$ is a primitive ray of its hexagon fan, and the
number of incidences at each base vertex $c$ is $b_c$.
\end{corollary}
\begin{proof}
Equality in the proof of Theorem~\ref{interact:bound} forces
$s=\ell=j+k-1$ and $j=\lfloor n/4\rfloor$. Connectedness therefore makes
the incidence graph a tree. Since the sum of all base capacities equals
the number of incidences, every base inequality is an equality, and
$\psi_i(t_{ic})=1$ for each nonzero twist. In a unimodular hexagon cone,
a nonzero integral vector with coefficient sum one is a primitive ray.
Finally $\sum_i\deg(i)=\ell$ gives the displayed excess-degree formula.
\end{proof}

The coefficient $5/4$ follows from the separate geometric restrictions
in~\eqref{interact:incidence-inequalities}. Each surface factor
contributes two dimensions and four Picard classes, but stability forces
at least two nonzero twists at that factor. Positivity of the base
primitive relations requires a total base dimension at least equal to
the number of such incidences. Thus the $j$ surface factors require
at least $2j$ base dimensions. Connectedness gives the remaining
inequality needed to bound the number of base factors.

The equality conditions are necessary, not sufficient: they do not
control the degrees of subspaces containing rays from several factors.
Section~\ref{sec:graphs} constructs the fans, and
Section~\ref{sec:subdivided} proves stability for independently
subdivided subcubic trees. Paths attain the bound when $4\mid n$;
the end attachments of Corollary~\ref{attach:sharp} give attainment in
every dimension $n\ge4$. Hence Theorem~\ref{interact:bound} is a sharp
bound in its stated bundle class. It is not an upper bound for arbitrary
smooth toric Fano varieties.
